\subsection{Compact groups of linear mappings}

Let E be a finite-dimensional vector space over R, C or H. Then End(E) is a finite-dimensional algebra over R, and the canonical topology on End(E) (§1, No.~10) is the topology of compact convergence. The group \(\operatorname{Aut}(E) = \mathbf{GL}(E)\) is an open subset of End(E), hence is a locally compact group. Let \((e_{1}, e_{2}, \ldots, e_{n})\) be a basis of E and, for every endomorphism u of E, let \(M(u) = (\alpha_{ij}(u))\) be the matrix of u with respect to this basis; to say that a subset S of End(E) is relatively compact in End(E) is equivalent to saying that the functions \(\alpha_{ij}(u)\) are bounded in S.

PROPOSITION 1. --- Let G be a subgroup of Aut(E). The following three properties are equivalent:

\begin{enumerate}
\def\labelenumi{(\roman{enumi})}
\item
  G is relatively compact in \(\operatorname{End}(\mathbf{E})\) ;
\item
  G is relatively compact in \(\operatorname{Aut}(\mathbf{E})\) ;
\item
  G leaves invariant a nondegenerate \(^{1}\) positive hermitian form on E. (iii) \(\Rightarrow\) (i): Suppose that G leaves invariant a nondegenerate positive hermitian form \(\Psi\) . Let \((e_{1},\ldots,e_{n})\) be an orthonormal basis for \(\Psi\) (Alg., Ch. IX, §6, No.~1, Cor. 1 of Th. 1). For every \(u\in G\) , let \((u_{ij})\) be its matrix with respect to \((e_{i})\) . For any j, we have \(\sum_{i=1}^{n}|u_{ij}|^{2}=1\) , thus \(|u_{ij}|\leqslant1\) for all i and j, which proves (i).
\item
  \(\Rightarrow\) (ii): This follows from GT, X, §3, No.~5, Cor. of Th. 4, taking into account the fact that the topology of \(\operatorname{End}(\mathbf{E})\) is that of compact convergence.
\item
  \(\Rightarrow\) (iii): Suppose that the closure \(\overline{\mathbf{G}}\) of \(\mathbf{G}\) in \(\operatorname{Aut}(\mathbf{E})\) is compact. Let \(\Phi\) be a nondegenerate positive hermitian form on \(\mathbf{E}\) . If the field of scalars is \(\mathbf{R}\) or \(\mathbf{C}\) , the giving of \(\Phi\) makes \(\mathbf{E}\) a finite-dimensional Hilbert space, and condition (iii) will result from the following lemma:
\end{enumerate}

Lemma 1. --- Let F be a Hilbert space, K a compact group, and \(s \mapsto U(s)\) a representation of K in the group of invertible elements of \(\mathcal{L}(\mathrm{F};\mathrm{F})\) , continuous for the topology of pointwise convergence. There exists a nondegenerate positive hermitian form \(\varphi\) on F such that

\[
\varphi (U (s) x, U (s) y) = \varphi (x, y)
\]

for all \(s \in \mathbf{K}\) , \(x \in \mathbf{F}\) , \(y \in \mathbf{F}\) , and such that the topological vector space structure of \(\mathbf{F}\) defined by \(\varphi\) (TVS, V, §1, No.~3) is identical to the original structure of \(\mathbf{F}\) .

Let \(\alpha\) be a Haar measure on \(\mathbf{K}\) . For any \(x, y\) in \(\mathbf{F}\) , the mapping \(s \mapsto (U(s)x|U(s)y)\) is continuous. Set

\[
\varphi (x, y) = \int (U (s) x | U (s) y) d \alpha (s).
\]

It is immediate that \(\varphi(x,y)\) is a sesquilinear form on \(\mathbf{F}\) . Since the set of endomorphisms \(U(s)\) is compact in \(\mathcal{L}_s(\mathrm{F};\mathrm{F})\) , there exists a constant \(\mathbf{M}\) such that \(\| U(s)\| \leqslant \mathbf{M}\) for all \(s\in \mathbf{K}\) . For every \(x\in \mathbf{F}\) , we therefore have

\[
\mathrm{M} ^ {- 1} \| x \| \leqslant \| U (s) x \| \leqslant \mathrm{M} \| x \|,
\]

whence the inequalities

\[
\mathbf {M} ^ {- 2} \alpha (\mathbf {K}) \| x \| ^ {2} \leqslant \varphi (x, x) \leqslant \mathbf {M} ^ {2} \alpha (\mathbf {K}) \| x \| ^ {2},
\]

which shows that \(\varphi\) is positive and nondegenerate, and that the norm \(\varphi(x,x)^{1/2}\) is equivalent to the norm \(\|x\|\) . Finally, for all \(t \in K\) ,

\[
\begin{array}{r l} & {\varphi (U (t) x, U (t) y) = \int (U (s t) x | U (s t) y) d \alpha (s)} \\ & {\qquad = \int (U (s) x | U (s) y) d \alpha (s) = \varphi (x, y).} \end{array}
\]

When the field of scalars is H, one argues exactly as before on replacing everywhere the function \(s \mapsto (U(s)x|U(s)y)\) by the function \(s \mapsto \Phi(sx, sy)\) defined on G, with values in H. This completes the proof of the proposition.

Remark. --- Let \(\Phi\) be a nondegenerate positive hermitian form on E. The unitary group \(\mathbf{U}(\Phi)\) is closed in \(\operatorname{Aut}(\mathbf{E})\) , hence is compact (Prop. 1). Prop. 1 also shows that every compact subgroup of \(\operatorname{Aut}(\mathbf{E})\) is contained in a subgroup of the form \(\mathbf{U}(\Phi)\) . If now \(\mathbf{U}(\Phi)\) is contained in a compact subgroup K of \(\operatorname{Aut}(\mathbf{E})\) , one sees that there exists a nondegenerate positive hermitian form \(\Phi'\) on E such that \(\mathbf{U}(\Phi) \subset \mathbf{K} \subset \mathbf{U}(\Phi')\) , and it follows easily (Exer. 1) that \(\Phi\) and \(\Phi'\) are proportional, whence \(\mathbf{U}(\Phi) = K\) . Thus the maximal compact subgroups of \(\operatorname{Aut}(E)\) are the subgroups of the form \(\mathbf{U}(\Phi)\) .

